\section*{Chapitre \ref{ch:qm:optimal-stopping-and-impulse-control} --- Arrêt optimal et contrôle impulsionnel}

\begin{solution}{exo:qm:optimal-stopping-and-impulse-control:1}
$b^* = (6 \times 300 \times 400/0.5)^{1/4} = 34.6$ millions~; $C^* = 2\sqrt{300 \times 400 \times 0.5/6} = 200$ dollars par jour~; $\sigma^2/b^{*2} = 1/3$ couverture
par jour, soit une fois tous les trois jours.
\end{solution}

\begin{solution}{exo:qm:optimal-stopping-and-impulse-control:2}
$U_2 = G_2 = (2, 0, 0)$ en $S_2 = 102, 100, 98$. $U_1 = \max(1, 1) = 1$ à 101 et $\max(0, 0) = 0$ à 99~; $U_0 = \max(0, 0.5) = 0.5$. À l'instant 0
continuer ($U > G$)~; à l'instant 1, l'arrêt est optimal (à 101, indifférence exacte~; à 99, rien n'est perdu).
\end{solution}

\begin{solution}{exo:qm:optimal-stopping-and-impulse-control:3}
Non~: le fait d'être le plus grand du mois n'est connu qu'en fin de mois, ce n'est donc pas un \omterm{def:qm:probability-at-speed:stopping}{temps d'arrêt}, et aucune politique ne peut
s'y couvrir.
\end{solution}

\begin{solution}{exo:qm:optimal-stopping-and-impulse-control:4}
En doublant $K$~: bande $\times 2^{1/4} = 1.19$ (41.2 millions), coût $\times\sqrt2$ (283 dollars). En doublant $\sigma$~: bande $\times\sqrt2$ (49.0 millions), coût $\times 2$
(400 dollars).
\end{solution}

\begin{solution}{exo:qm:optimal-stopping-and-impulse-control:5}
$b^* = (3 \times 25 \times 400/2)^{1/3} = 24.7$ millions~; coût $25 \times 400/49.3 + 0.5 \times 24.7^2/3 = 304$ dollars par jour.
\end{solution}

\begin{solution}{exo:qm:optimal-stopping-and-impulse-control:6}
En deçà du seuil $V(x) = Ae^{\theta x}$ (la solution de $\mathcal LV = rV$ qui reste bornée quand $x \to -\infty$)~; raccordement des valeurs $Ae^{\theta b} = b - c$ et
\omterm{def:qm:optimal-stopping-and-impulse-control:fbp}{recollement C1} (smooth pasting) $A\theta e^{\theta b} = 1$ donnent $\theta(b - c) = 1$. Avec les chiffres, $\theta = (-0.5 + \sqrt{0.25 + 0.16})/4 = 0.0351$ et $b^* = 5 + 28.5 =
33.5$ pb.
\end{solution}

\begin{solution}{exo:qm:optimal-stopping-and-impulse-control:7}
197.5 dollars par jour pour 390 contrôles et 199.6 pour 1\,560, contre 200. Avec des contrôles discrets, l'exposition dépasse la bande avant une
couverture, de sorte que les cycles durent un peu plus longtemps et qu'il y a moins d'ordres~; le surcroît de charge de risque lié au dépassement est inférieur aux frais économisés, et l'écart se referme
comme $\sqrt{\Delta t}$.
\end{solution}

\begin{solution}{exo:qm:optimal-stopping-and-impulse-control:8}
Pour $b = 10$, la politique coûte $300 \times 400/100 + 0.5 \times 100/6 = 1\,208$ dollars par jour, soit six fois l'optimum~: les frais d'une douzaine
de couvertures par jour dépassent largement le risque qu'elles suppriment. Le risque et les frais doivent être arbitrés~; réduire le risque n'est pas gratuit.
\end{solution}

\begin{solution}{pb:qm:optimal-stopping-and-impulse-control:1}
\textbf{1.} $C(b) = K\sigma^2/b^2 + \gamma b^2/6$.
\textbf{2.} $b^* = (6K\sigma^2/\gamma)^{1/4} = 34.6$ millions.
\textbf{3.} 200 dollars par jour.
\textbf{4.} Tous les trois jours en moyenne.
\textbf{5.} $300 + 0.5 \times 400/2 = 400$ dollars par jour~: la bande fait économiser 50~\%.
\textbf{6.} 29.1 millions (16~\% plus étroite) et 141 dollars par jour.
\textbf{7.} 49.0 millions et 400 dollars par jour.
\textbf{8.} La moitié~: à l'optimum, les frais et la charge de risque sont égaux.
\textbf{9.} 197.5 et 199.6 dollars par jour.
\textbf{10.} Des contrôles discrets laissent l'exposition dépasser la bande, allongeant les cycles et réduisant le nombre d'ordres.
\textbf{11.} 24.7 millions et 304 dollars par jour.
\textbf{12.} Déclenchement à 42.4 millions, retour à 12.8 millions, 418 dollars par jour.
\textbf{13.} Chaque million échangé coûte 25 dollars, de sorte qu'arrêter avant zéro économise un coût proportionnel, et le cycle suivant, qui part
plus près de la bande, reste assez long pour amortir les frais fixes.
\textbf{14.} Le \omterm{def:qm:optimal-stopping-and-impulse-control:singular}{contrôle singulier} échange de façon infinitésimale, seulement lorsque l'exposition touche la bande, avec un montant total échangé fini~; le \omterm{def:qm:optimal-stopping-and-impulse-control:impulse}{contrôle impulsionnel}
effectue un nombre fini d'échanges discrets, chacun payant les frais fixes.
\textbf{15.} $b^* = 33.5$ pb sous forme close~; la grille quotidienne s'arrête à 32.2 pb et valorise la position à zéro à 8.79 pb contre 8.80.
\textbf{16.} Comme charge de variance, oui~: la variance du P\&L d'une exposition ouverte est proportionnelle à $X^2$. Une charge proportionnelle à $|X|$ (une
limite de type VaR) modifie les exposants des lois de la bande.
\textbf{17.} Plus large~: des flux qui se compensent rendent l'attente moins coûteuse, puisqu'une partie de l'exposition disparaît sans échange.
\textbf{18.} Des expositions corrélées se compensent en partie~; fixer une bande sur le risque net conduit à couvrir moins souvent que de fixer une bande pour chaque devise.
\textbf{19.} Résultat nommé~: \emph{bande de couverture}~: la bande optimale de non-transaction vaut $\pm(6K\sigma^2/\gamma)^{1/4} = \pm 34.6$ millions, coûtant 200 dollars par jour,
soit la moitié du coût de la couverture quotidienne~; la bande croît comme la racine quatrième des frais fixes par ordre (16~\% plus étroite si ces frais sont divisés par deux).
\textbf{20.} À l'intérieur de la bande, le risque supprimé par une couverture vaut moins que les frais qu'elle coûte, et attendre conserve la possibilité de couvrir une
exposition plus importante pour les mêmes frais.
\end{solution}

\begin{solution}{iq:qm:optimal-stopping-and-impulse-control:1}
Une bande~: ne rien faire tant que l'exposition reste dans $\pm b$, revenir à zéro au bord, avec $b = (6K\sigma^2/\gamma)^{1/4}$ pour des frais fixes $K$, une volatilité des flux
$\sigma$ et une charge de risque $\gamma X^2$. Elle égalise les frais et la charge de risque.
\iqlookfor{\omterm{def:qm:optimal-stopping-and-impulse-control:notrade}{région de non-transaction}, le compromis, et les lois d'échelle.}
\end{solution}

\begin{solution}{iq:qm:optimal-stopping-and-impulse-control:2}
À la frontière optimale, la valeur de l'attente et celle de l'arrêt coïncident (raccordement des valeurs), de même que leurs pentes. Si la valeur de l'attente
rejoignait le gain sous un angle, déplacer légèrement la frontière ferait monter la valeur~: le seuil ne serait pas optimal.
\iqlookfor{optimalité du premier ordre de la frontière.}
\end{solution}

\begin{solution}{iq:qm:optimal-stopping-and-impulse-control:3}
Avec des frais fixes, le coût par unité de temps est égal aux frais multipliés par la fréquence de couverture $\sigma^2/b^2$, plus une charge de risque $\propto b^2$~: $b^4 \propto K$. Avec un
coût proportionnel, les échanges ont lieu à la frontière, à un rythme de $\sigma^2/b$ en volume, de sorte que le coût vaut $c\sigma^2/b$ plus un terme $\propto b^2$~: $b^3 \propto c$.
\iqlookfor{à quelle fréquence et en quelle quantité on échange en fonction de $b$.}
\end{solution}

\begin{solution}{iq:qm:optimal-stopping-and-impulse-control:4}
La plus petite \omterm{def:qm:probability-at-speed:martingale}{surmartingale} dominant le processus de gain, calculée à rebours comme le maximum entre arrêter maintenant et la valeur espérée
de poursuivre~; le premier instant où elle égale le gain est optimal. Le prix d'une option américaine est l'\omterm{def:qm:optimal-stopping-and-impulse-control:snell}{enveloppe de Snell} de son gain actualisé sous la
\omterm{def:qm:girsanov-and-changes-of-numeraire:emm}{probabilité risque-neutre} (livre~5, chapitre~6).
\iqlookfor{la récurrence à rebours et son optimalité.}
\end{solution}

\begin{solution}{iq:qm:optimal-stopping-and-impulse-control:5}
Échanger jusqu'à un niveau de remise strictement à l'intérieur de la bande, jamais sur la bande elle-même, afin qu'une couverture ne déclenche pas aussitôt une autre~;
mesurer l'exposition nette des couvertures en attente~; lisser les flux d'exposition bruités~; plafonner le nombre de couvertures par intervalle~; consigner chaque décision avec l'exposition et la bande en vigueur~;
calculer la bande à partir des données courantes, et non d'une constante.
\iqlookfor{hystérésis, ordres en cours, et traçabilité.}
\end{solution}

\begin{solution}{iq:qm:optimal-stopping-and-impulse-control:6}
$f(x) = \E_x\int_0^\tau X^2\,dt$ est solution de $\tfrac12\sigma^2f^{\prime\prime} = -x^2$ avec $f(\pm b) = 0$, donc $f(x) = (b^4 - x^4)/(6\sigma^2)$ et $f(0) = b^4/(6\sigma^2)$.
\iqlookfor{formule de Dynkin et problème aux limites.}
\end{solution}
